\chapter{Glossario discorsivo dei metodi}
\label{ch:glossario_discorsivo}

Questo capitolo spiega \emph{cos'\`e} ogni metodo che $\pi$Tantum
usa internamente, con un linguaggio che si possa leggere
``dal divano''. Per ogni metodo trovi quattro cose in
sequenza: l'idea, come funziona, quando usarlo, un esempio
legato al timetabling. Niente formule a meno che non aggiungano
chiarezza.

\section{CP-SAT (Constraint Programming over SAT)}

\textbf{Cos'\`e e perch\'e esiste.} CP-SAT \`e il motore con cui
$\pi$Tantum trova le soluzioni. \`E un \emph{constraint
solver}: gli dai un mucchio di variabili (per esempio: ``in
quale slot va la lezione di mate della 1A?''), un mucchio di
vincoli (``due lezioni della stessa classe non possono stare
nello stesso slot''; ``il prof Borghi non insegna sabato''),
e gli chiedi di trovare un assegnamento di valori che li
rispetti tutti, possibilmente ottimizzando un criterio.

A differenza dei solver \emph{lineari} (LP/MIP) che lavorano
su disuguaglianze numeriche, CP-SAT lavora su variabili
booleane e intere, e sa esprimere vincoli "logici" complessi
(disgiunzioni, implicazioni, conteggi) in modo nativo. Questo
lo rende particolarmente adatto al timetabling, dove la
maggior parte dei vincoli sono di natura combinatoria, non
geometrica.

\textbf{Come funziona.} CP-SAT alterna due fasi:
\begin{itemize}
\item \textbf{Propagazione}: data una scelta parziale (es. ``ho
  deciso che mate-3A va in luned\`i 8\textsuperscript{a}''),
  il solver \emph{deduce} automaticamente quello che ne
  consegue (``allora ita-3A non pu\`o andare in luned\`i
  8\textsuperscript{a}''; ``allora il prof Borghi \`e occupato
  in quello slot'', ecc.). Questo restringe lo spazio delle
  scelte rimanenti senza dover provare tutte le combinazioni.
\item \textbf{Ricerca}: quando la propagazione esaurisce le
  conseguenze automatiche, il solver fa un \emph{salto
  d'azzardo} (es. ``provo a mettere ita-3B in marted\`i 9'')
  e si rimette a propagare. Se a un certo punto trova una
  contraddizione, fa \emph{backtracking}: torna indietro,
  marca quella scelta come ``no'', e prova un'alternativa.
  Impara dai fallimenti aggiungendo \emph{clausole imparate}
  che evitano di ripetere lo stesso errore.
\end{itemize}

\textbf{Quando usarlo.} CP-SAT \`e il default di $\pi$Tantum
per le Phase A e B. Per scuole di dimensioni normali (fino a
$\sim$80 classi) basta da solo. Per istanze pi\`u grandi
serve combinarlo con la decomposizione spettrale o con
metaeuristiche.

\textbf{Esempio.} Tipica iterazione: dopo aver propagato per
qualche secondo, il solver ha gi\`a piazzato il 70\% delle
lezioni del triennio (perch\'e sono molto vincolate dai lab e
dalle compresenze) ed esplora le combinazioni rimanenti per
le classi del biennio (pi\`u flessibili). Trova una
contraddizione (un docente avrebbe due classi nello stesso
slot), backtracking, ne prova un'altra. Tipicamente in 1-3
minuti chiude su una soluzione feasible.

\section{Decomposizione spettrale}

\textbf{Cos'\`e e perch\'e esiste.} Quando la scuola \`e grande
(>120 classi), il problema diventa troppo grosso perch\'e
CP-SAT lo digerisca in tempi ragionevoli. La decomposizione
spettrale risolve questo dividendo la scuola in piccoli
``cluster'' di classi che condividono pochi docenti. Ogni
cluster diventa un sotto-problema piccolo, risolvibile in
parallelo.

L'idea \`e: i conflitti veri (due docenti che competono per
lo stesso slot) avvengono fra classi che condividono lo
stesso docente. Se due classi non hanno docenti in comune,
si possono pianificare totalmente in parallelo.

\textbf{Come funziona.} Si costruisce il grafo delle classi:
nodi = classi, peso dell'arco fra due classi = numero di
docenti condivisi. Si calcolano gli autovettori del
\emph{Laplaciano} di quel grafo (una matrice che misura,
per ogni nodo, quanto si ``differenzia'' dai vicini). Gli
autovettori associati ai pi\`u piccoli autovalori
diversi da zero indicano \emph{direzioni di separazione}
naturali: progettando i nodi su quelle direzioni, si vedono
chiaramente i cluster.

Un colpo di k-means su questa proiezione restituisce la
partizione finale.

\textbf{Quando usarlo.} Default ON in $\pi$Tantum quando la
scuola ha $\geq 120$ classi. Per scuole pi\`u piccole non
porta vantaggi e aggiunge overhead.

\textbf{Esempio.} Per una scuola di 250 classi, una
decomposizione in 4 cluster produce sotto-problemi di 60-70
classi ciascuno, risolvibili in parallelo. I docenti-ponte
(quelli che insegnano in classi di pi\`u cluster) sono il
``vincolo di ricucitura'' che la fase finale risolve come
problema piccolo a se' stante.

\section{LNS (Large Neighborhood Search)}

\textbf{Cos'\`e e perch\'e esiste.} Quando hai gi\`a una
soluzione (anche sub-ottima), spesso \`e pi\`u rapido
\emph{aggiustarla} che ricominciare. LNS prende la soluzione
attuale, ne ``distrugge'' una porzione (es. tutte le lezioni
di un certo giorno) e la ricostruisce con CP-SAT cercando un
miglioramento. Iterando, esplora tante varianti vicine alla
soluzione corrente.

\textbf{Come funziona.} A ogni iterazione: scegli un
``vicinato'' (es. day = marted\`i), libera tutte le variabili
in quel vicinato, lancia CP-SAT con un piccolo budget di
tempo per re-piazzarle ottimizzando lo SOFT score. Se il
risultato \`e migliore, lo accetti.

\textbf{Quando usarlo.} Default ON nella pipeline integrata,
dopo Phase B. Adatto al ``polishing'' di una soluzione
gi\`a feasible.

\textbf{Esempio.} Dopo Phase B la soluzione ha 22 buchi nei
docenti. LNS cicla: prova a re-fare luned\`i (scende a 20
buchi), marted\`i (scende a 19), mercoled\`i (resta 19),
\dots, dopo qualche minuto si stabilizza intorno a 14-15.

\section{ALNS (Adaptive LNS)}

\textbf{Cos'\`e e perch\'e esiste.} LNS classico ha un
problema: usa sempre lo stesso tipo di vicinato (es. ``un
giorno''). Su alcune istanze \`e perfetto, su altre \`e
sterile. ALNS lo evolve: invece di una sola scelta di
vicinato, ne ha sei (un giorno intero, un cluster di classi,
le ore peggiori della settimana, una giornata di un docente,
una giornata di un'aula, una settimana di una sola classe) e
tre modi di ricostruire (CP-SAT mini-window, greedy SOFT,
BFS riempimento). A ogni iterazione sceglie quale combinazione
usare con probabilit\`a proporzionale a un \emph{punteggio}
che cresce per le combinazioni storicamente vincenti e
decresce per le altre.

\textbf{Come funziona.} \`E un \emph{roulette wheel selector}
sopra punteggi esponenzialmente decadenti. Ogni operatore
parte con punteggio neutro; ogni volta che produce un
miglioramento, il punteggio sale; ogni volta che fallisce,
scende. Aggiunge anche un'accettazione \emph{simulated
annealing}-like: occasionalmente accetta peggioramenti per
sfuggire ai minimi locali.

\textbf{Quando usarlo.} Sostituisce o segue il LNS classico
nella pipeline. Default ON.

\textbf{Esempio.} Su una scuola con cluster ben separati,
l'ALNS impara dopo poche iterazioni che ``cluster di classi''
porta i miglioramenti pi\`u grossi e ``giornata di un'aula''
quasi mai funziona. Lo dice nei log finali con i punteggi.

\section{SA (Simulated Annealing)}

\textbf{Cos'\`e e perch\'e esiste.} Il nome viene dalla
metallurgia: come un metallo riscaldato e poi raffreddato
lentamente forma una struttura cristallina pi\`u stabile,
l'algoritmo accetta inizialmente anche peggioramenti per
esplorare lo spazio delle soluzioni, e con il tempo (la
``temperatura'' che scende) diventa pi\`u selettivo,
accettando solo miglioramenti.

\textbf{Come funziona.} A ogni iterazione genera una mossa
random (es. scambia due lezioni). Se migliora lo SOFT, accetta;
se peggiora di $\Delta$, accetta con probabilit\`a
$e^{-\Delta/T}$ dove $T$ \`e la temperatura corrente. La
temperatura inizia alta (es. $T_0 = 10$) e si moltiplica per
$\alpha < 1$ (es. $\alpha = 0.995$) a ogni iterazione, quindi
decresce esponenzialmente.

\textbf{Quando usarlo.} \`E perfetto quando si sospetta di
essere in un minimo locale. Default ON dopo LNS/ALNS.

\textbf{Esempio.} Dopo LNS ti sei stabilizzato a SOFT=420.
Lanci SA con T0=10, alpha=0.995, budget 60s. Nei primi
secondi, $T$ \`e alto e accetta peggioramenti fino a 460-470;
poi $T$ scende e SA si concentra su miglioramenti, finendo
per esempio a 405.

\section{TS (Tabu Search)}

\textbf{Cos'\`e e perch\'e esiste.} Risolve un problema della
ricerca locale: i loop. Se l'unica mossa migliorante
\`e ``A $\to$ B'' e poi l'unica migliorante da B \`e ``B
$\to$ A'', sei intrappolato. TS evita questo tenendo una
\emph{lista tabu} delle ultime mosse fatte: per un certo
numero di iterazioni, quelle mosse sono vietate. L'algoritmo
\`e cos\`i forzato a esplorare zone diverse.

\textbf{Come funziona.} A ogni iterazione, valuta tutte le
mosse possibili \emph{tranne} quelle in tabu. Sceglie la
migliore, la fa, aggiunge l'inversa alla lista tabu (con un
``tempo di scadenza''). La lista \`e di dimensione fissa
(parametro \texttt{tabu\_size}, default 80).

\textbf{Quando usarlo.} Default ON dopo SA. Particolarmente
efficace su istanze con plateau (zone dove molti vicini hanno
lo stesso valore).

\textbf{Esempio.} Senza TS, scambiare due lezioni $L_1$ e
$L_2$ produce SOFT=400; ri-scambiarle riproduce SOFT=405.
SA potrebbe oscillare. TS dopo lo scambio mette ``inverso di
$L_1 \leftrightarrow L_2$'' in tabu, quindi al passo
successivo prover\`a $L_3 \leftrightarrow L_4$, esplorando
una direzione nuova.

\section{ILS (Iterated Local Search)}

\textbf{Cos'\`e e perch\'e esiste.} Combina ricerca locale
con perturbazioni periodiche. Idea: dopo che la ricerca
locale si \`e stabilizzata su un minimo locale, scuoti la
soluzione (perturbazione: tante mosse random) e fai ricerca
locale da capo. Iterando, esplori bacini di attrazione
diversi.

\textbf{Come funziona.} Loop di ``cycles'' (default 3): in
ognuno fai ricerca locale (TS) per un budget; se sei migliorato,
salvi; poi perturba (rimuovi e ri-piazza un cluster random
di classi); ricomincia.

\textbf{Quando usarlo.} Default ON come ultima fase prima
del rooms assignment. \`E il \emph{closer} della pipeline
metaeuristica.

\textbf{Esempio.} Cycle 1 chiude su SOFT=400; perturba e
ricomincia da SOFT=480; cycle 2 chiude su SOFT=395; perturba
e ricomincia da SOFT=470; cycle 3 chiude su SOFT=392.

\section{VNS (Variable Neighborhood Search)}

\textbf{Cos'\`e e perch\'e esiste.} VNS \`e una rifinitura
sistematica. Cicla attraverso vicinati di dimensione
crescente: prima prova 1-swap (scambia due lezioni), poi
2-swap (due paia), poi 3-chain (catena di 3 mosse), poi k-opt
(catena di 4-6). Ogni volta che trova un miglioramento riparte
dal vicinato 1; quando un intero ciclo passa senza
miglioramenti, si ferma -- significa che ha raggiunto un
ottimo locale rispetto a tutti i vicinati.

\textbf{Come funziona.} For ogni vicinato, esplora fino a un
budget di iterazioni; al primo miglioramento accetta e torna
al vicinato 1. Quando il giro $1 \to 4$ produce zero
miglioramenti, esce.

\textbf{Quando usarlo.} Come rifinitura finale, dopo
LNS/ALNS/SA/TS. Default OFF (lo accendi quando vuoi spremere
il massimo).

\textbf{Esempio.} Hai SOFT=395 dopo ILS. VNS cicla: 1-swap
trova un miglioramento a 392; riparte. 1-swap trova ancora a
389; riparte. 1-swap zero; 2-swap trova a 385; riparte. Dopo
qualche minuto chiude a 378.

\section{Hall's theorem pre-check}

\textbf{Cos'\`e e perch\'e esiste.} Prima di lanciare un
solver lungo, conviene controllare se il problema \`e
\emph{strutturalmente} risolvibile. Il teorema di Hall (1935)
dice: in un \emph{matching bipartito} (assegnare ognuno di N
elementi a uno di M elementi compatibili), una soluzione
esiste sse per ogni sottoinsieme $S$ degli N elementi, il
loro vicinato (gli M compatibili con almeno uno) ha
cardinalit\`a $|N(S)| \geq |S|$.

Nel timetabling l'analogo pesato \`e: per ogni sottoinsieme di
docenti, la loro capacit\`a totale di ore deve coprire la
domanda totale di ore delle materie che insegnano. Se una
materia richiede 40 ore/settimana ma i docenti qualificati ne
hanno solo 30 totali, il modello \`e gi\`a infeasible: non
c'\`e bisogno di lanciare CP-SAT.

\textbf{Come funziona.} Tre controlli in cascata: per ogni
materia controlla che esista almeno un docente qualificato;
per ogni materia controlla che la capacit\`a totale dei
docenti qualificati copra la domanda; campiona N sottoinsiemi
random di docenti e verifica che la domanda delle materie
\emph{esclusive} a ciascun sottoinsieme sia coperta dalla
capacit\`a del sottoinsieme.

\textbf{Quando usarlo.} Sempre prima di Phase A. Default ON
nella pipeline integrata; se trova violazioni, interrompe
immediatamente con un report.

\textbf{Esempio.} Hai aggiunto in mezzo all'anno una nuova
materia ``Diritto'' al triennio (3 ore/settimana per 9 classi
= 27 ore) ma hai un solo docente abilitato a Diritto con
massimo 18 ore. Hall te lo dice in 100 millisecondi: ``Diritto:
domanda 27h > capacit\`a docenti 18h''. Risparmi 10 minuti di
solver lungo.

\section{Lagrangian Relaxation con sub-gradient}

\textbf{Cos'\`e e perch\'e esiste.} Quando un problema ha
``vincoli scomodi'' (vincoli che spezzano una struttura
altrimenti separabile), una tecnica classica \`e
\emph{rilassarli}: li rimuovi dal modello e li reintroduci
come penalit\`a moltiplicate per coefficienti chiamati
\emph{moltiplicatori di Lagrange} ($\lambda$). Iterando, i
$\lambda$ si aggiustano per spingere la soluzione del modello
rilassato a rispettare comunque i vincoli originari. \`E un
modo di trasformare un problema ``a blocchi accoppiati'' in
una sequenza di sotto-problemi indipendenti.

Nel timetabling questo \`e utile dopo la decomposizione
spettrale: i cluster sarebbero indipendenti se non fosse per
i \emph{docenti-ponte} (docenti che insegnano in pi\`u
cluster). Lagrangian relaxation rilassa la non-sovrapposizione
dei docenti-ponte, e itera per riconciliare le scelte dei
cluster.

\textbf{Come funziona.} Il \emph{subgradient method} aggiorna
$\lambda$ a ogni iterazione: $\lambda_{k+1} = \max(0, \lambda_k +
\alpha_k g_k)$, dove $g_k$ \`e la violazione del vincolo
rilassato e $\alpha_k = \alpha_0/(1+k)$ \`e una sequenza di
``passi'' che diminuisce a passi piccoli, garantendo
convergenza al limite (in teoria) ma in pratica producendo
soluzioni utili in poche iterazioni.

\textbf{Quando usarlo.} Scuole con cluster ben separati ma
con docenti-ponte fastidiosi. Default OFF (avanzato).

\textbf{Esempio.} Cluster A (40 classi), Cluster B (40
classi), 5 docenti-ponte. Senza Lagrangian, A e B vanno
risolti insieme: 80 classi. Con Lagrangian, vengono
risolti separatamente in $\sim$3 iterazioni: ogni cluster
diventa 40 classi (molto pi\`u veloce), e $\lambda$ converge
a valori che impediscono ai docenti-ponte di sovrapporsi.

\section{Column Generation (Dantzig-Wolfe)}

\textbf{Cos'\`e e perch\'e esiste.} Quando un problema \`e
enorme, conviene non enumerarlo tutto: ne generi solo le
parti che servono, on-demand. Column Generation \`e questa
filosofia applicata alla programmazione lineare: hai un
problema ``master'' che sceglie da un catalogo di colonne
(soluzioni candidate); a ogni iterazione un sotto-problema
genera nuove colonne ``promettenti'' (con \emph{reduced cost}
negativo) e le aggiunge al catalogo.

Nel timetabling: ogni docente ha un catalogo di possibili
``settimane-tipo'' (orari completi per quel docente, gi\`a
feasible localmente). Il master sceglie quale settimana-tipo
prendere per ognuno; il sotto-problema (uno per docente) genera
nuove settimane-tipo guidate dal duale del master.

\textbf{Come funziona.} Loop: (1) il master sceglie una
combinazione corrente di colonne; (2) calcola i prezzi duali
sui vincoli di copertura; (3) ogni sotto-problema produce la
nuova colonna pi\`u attraente data quei prezzi; (4) se ce
n'\`e una con reduced cost negativo, aggiungila e ripeti;
altrimenti hai trovato l'ottimo del rilassamento LP.

\textbf{Quando usarlo.} Solo per scuole molto grandi
(>200 classi). Default OFF.

\textbf{Esempio.} 250 classi, 180 docenti. Enumerare tutte
le settimane-tipo possibili sarebbe astronomicamente troppo;
Column Generation parte con 3 settimane-tipo per docente
(540 colonne), ne aggiunge $\sim$50 per iterazione, e
converge in 5-10 iterazioni a un orario ottimale rilassato.

\section{Monte Carlo Sensitivity}

\textbf{Cos'\`e e perch\'e esiste.} Quando hai una soluzione,
ti chiedi: \emph{\`e gi\`a vicina al meglio possibile, o c'\`e
ancora margine}? La sensitivity analysis Monte Carlo risponde
generando N varianti random della soluzione (ognuna ottenuta
applicando 30 mosse atomiche random feasible) e misurando il
loro SOFT score. Se la varianza \`e piccola, sei vicino a un
minimo locale; se \`e grande, hai ancora margine.

\textbf{Come funziona.} Genera N (default 100) walk random
di 30 mosse ciascuno, accettando solo quelle che restano
HARD-feasible. Per ognuno calcola SOFT. Riporta media,
deviazione standard, percentili, coefficiente di variazione
$CV = \sigma/\mu$.

\textbf{Quando usarlo.} Dopo aver lanciato la pipeline, per
decidere se vale la pena spendere altro tempo in metaeuristiche.
$CV < 0.05$ = stai gi\`a sull'ottimo locale; $CV > 0.15$ =
c'\`e potenziale di miglioramento.

\textbf{Esempio.} La tua soluzione corrente ha SOFT=400. Monte
Carlo restituisce media=402, std=8, $CV \approx 0.02$. Conclusione:
la pipeline ha gi\`a trovato un buon minimo locale; nuovi
giri non porteranno miglioramenti significativi.

\section{Modularit\`a, betweenness, densit\`a (analisi del bipartito)}

\textbf{Cos'\`e e perch\'e esiste.} Tre metriche dalla teoria
dei grafi che ti dicono quanto la struttura della tua scuola
sia ``decomponibile''. Sono interessanti perch\'e predicono
quanto certe tecniche (decomposizione spettrale, parallelismo)
funzioneranno bene.

\textbf{Cosa significano in pratica.}
\begin{itemize}
\item \textbf{Modularit\`a (Newman-Girvan)}: misura quanto
  una partizione \emph{gi\`a calcolata} \`e ``naturale''.
  Valore in $[-1, 1]$. $> 0.4$ = cluster molto netti, la
  decomposizione spettrale \`e perfetta. $< 0.2$ = la scuola
  \`e troppo intrecciata, la decomposizione produrr\`a
  cluster artificiosi.
\item \textbf{Betweenness centrality}: per ogni nodo, conta
  in quanti percorsi minimi del grafo passa. I nodi ad alta
  betweenness sono i ``ponti'': se hanno SOFT alto, fanno
  collassare la qualit\`a globale.
\item \textbf{Densit\`a}: $\frac{2 |E|}{|V|(|V|-1)}$, frazione
  di archi presenti rispetto al massimo. Bassa densit\`a
  ($< 0.1$) = scuola sparsa, decomposizione ottima. Alta
  densit\`a ($> 0.5$) = quasi cricca, decomposizione
  inutile.
\end{itemize}

\textbf{Quando usarli.} Diagnostici da consultare prima di
investire in metaeuristiche pesanti: ti dicono che strategia
ha senso.

\textbf{Esempio.} Liceo medio: modularit\`a 0.43, betweenness
top = ``Lingua straniera'' (5 prof condivisi su tutti gli
indirizzi), densit\`a 0.18. Diagnosi: la decomposizione
spettrale produrr\`a 4 cluster netti (Liceo Classico, Liceo
Scientifico, Liceo Linguistico, ITIS); i prof di lingua sono
i ``ponti'' da gestire con cura (Lagrangian relaxation aiuta).

\section{Distribuzioni e istogrammi}

\textbf{Cosa misurano.} 5 viste statistiche dell'orario
prodotto: distribuzione del carico orario dei docenti,
distribuzione del carico delle classi per giorno, heatmap
materia $\times$ slot, occupazione per slot, test di
goodness-of-fit (KS contro uniforme, chi-quadro per
giornate).

\textbf{Cosa significano in pratica.}
\begin{itemize}
\item \emph{Carico docenti}: se la distribuzione \`e
  schiacciata su un numero (es. tutti 18) il sistema sta
  rispettando i tetti di cattedra. Se ha una coda lunga (un
  prof con 22, un altro con 12), c'\`e sbilanciamento.
\item \emph{Heatmap materia $\times$ slot}: zone scure dove
  cadono troppe ore di una stessa materia in slot adiacenti.
\item \emph{KS-test contro uniforme}: $p > 0.05$ = i carichi
  sono compatibili con una distribuzione uniforme (giusto);
  $p < 0.05$ = c'\`e una distribuzione non casuale (alcuni
  prof troppo carichi).
\end{itemize}

\textbf{Quando usarli.} Per dare un giudizio qualitativo
sull'orario prodotto, oltre al solo SOFT score numerico.

\section{Correlazioni e regressioni}

\textbf{Cosa misurano.} Tre regressioni statistiche
sull'orario corrente:
\begin{enumerate}
\item OLS: ore di carico docente $\to$ numero di buchi.
\item OLS: numero di materie diverse del docente $\to$
  contribuzione SOFT.
\item Logit: saturazione classe $\to$ probabilit\`a di sesta
  ora.
\end{enumerate}

Per ognuna ottieni coefficienti, p-value, R$^2$.

\textbf{Cosa significano in pratica.} ``I docenti con carico
> 20 ore hanno in media 2 buchi in pi\`u ($p < 0.001$,
$R^2 = 0.6$)'' = c'\`e un effetto sistematico documentabile.
``La saturazione classe non predice la 6\textsuperscript{a}
ora ($p = 0.4$)'' = il sistema sta distribuendo bene le ore
e non concentra le classi sature in 6\textsuperscript{a}.

\textbf{Quando usarle.} Per riportare al collegio docenti
dati \emph{quantitativi} sulla qualit\`a dell'orario, anzich\'e
solo affermazioni qualitative.

\section{DSL generico}

\textbf{Cos'\`e e perch\'e esiste.} I sub-DSL precedenti
(query lista, logical DNF, objective Phase A, introspettivo)
risolvevano ognuno un caso, ma erano grammatiche separate.
Ogni fix andava replicato 4 volte. Il \emph{DSL generico}
unifica: una grammatica sola, parser unico, quattro back-end di
compilazione.

\textbf{Come si rapporta al modello dati.} Le sorgenti
iterabili (\texttt{lessons}, \texttt{teachers},
\texttt{classes}, \texttt{classrooms}, \texttt{students},
\texttt{groups}, \dots) sono viste pre-calcolate del modello
SQLAlchemy: a ogni valutazione, il sistema ``materializza''
queste viste in dict Python che il parser interroga.

\textbf{Come si traduce.} Un \texttt{forall l in lessons
where ...: predicate} si traduce, per il backend EVAL, in un
ciclo Python su tutti gli elementi che soddisfano il filtro,
applicando il predicato. Per il backend LOGIC (vincoli
logical DNF) si traduce in clausole disgiuntive di
letterali. Per il backend OBJECTIVE si traduce in una somma
pesata aggiunta alla funzione obiettivo della Phase A. Stessa
sintassi, traduzioni diverse.

\textbf{Quando usarlo.} Per qualunque vincolo che non rientri
nelle preferenze standard. Vedi capitolo dedicato per la
gallery di 30+ esempi.

\section{Stati 5-colori delle matrici}

\textbf{Cos'\`e.} Le matrici di disponibilit\`a (docente,
classe, aula) hanno 5 stati ortogonali per ogni cella, con
colori distintivi per leggerli a colpo d'occhio:

\begin{itemize}
\item \textbf{HARD} (rosso): cella vietata in modo assoluto.
\item \textbf{ENFORCED} (rosso scuro): cella che \emph{deve}
  esserci occupata.
\item \textbf{PREFERRED} (verde): preferenza positiva (bonus).
\item \textbf{DISLIKED/SOFT} (giallo): preferenza negativa
  (penalit\`a).
\item \textbf{ALLOWED} (bianco): neutro.
\end{itemize}

\textbf{Scorciatoie tastiera.} Selezionata una cella o un
range, premere \texttt{H/E/P/D/A} la imposta in HARD/ENFORCED/
PREFERRED/DISLIKED/ALLOWED. \texttt{N} per ``neutro'' (alias
ALLOWED). Click + drag per selezionare un range; doppio click
sull'header di colonna o riga per applicare a tutta la
colonna/riga.

\section{Tag aule e tag studenti}

\textbf{Cos'\`e.} Etichette libere che il coordinatore attacca
a un'aula o a uno studente. Sono \emph{many-to-many}: un'aula
pu\`o avere tanti tag, un tag pu\`o essere su tante aule.

\textbf{Perch\'e esistono.} I tag rendono i vincoli pi\`u
robusti al cambiamento dell'anagrafica. Invece di scrivere
``solo queste tre aule precise'' (lista che andrebbe
aggiornata a ogni modifica), scrivi ``aule taggate
matematica'': se domani aggiungi una nuova aula taggata
allo stesso modo, il vincolo la include automaticamente.

\textbf{Casi d'uso tipici.}
\begin{itemize}
\item Tag aule: \texttt{matematica}, \texttt{proiettore},
  \texttt{lab\_lingue}, \texttt{piano\_terra},
  \texttt{accessibile\_disabili}.
\item Tag studenti: \texttt{BES} (Bisogni Educativi Speciali),
  \texttt{DSA}, \texttt{debito\_matematica\_4},
  \texttt{pcto\_Acme}, \texttt{studente\_atleta}.
\end{itemize}

I tag NON sostituiscono i gruppi articolati (che sono unit\`a
operative di scheduling); sono metadato passivo per filtrare
e raggruppare nell'interfaccia.

\section{Run telemetry e visualizzazione live}

\textbf{Cos'\`e.} Ogni esecuzione del solver produce una
\emph{serie temporale} di sample: ad ogni iterazione il
sistema annota timestamp, fase corrente, valore obiettivo,
mosse accettate/rifiutate, vincoli HARD violati. Tutto
salvato in DB e accessibile dal detail page del run.

\textbf{Cosa vedi.} Un grafico interattivo (Apache ECharts)
con l'evoluzione del SOFT score nel tempo, una linea diversa
per ogni stage (LNS in azzurro, SA in verde, TS in giallo,
\dots), markers verticali sui cambi di stage. Statistiche
aggregate per stage in tabella (n. iterazioni, durata, best
obj). Esportazione in JSON per analisi esterne.

\textbf{Live mode.} Per i run in corso, la pagina poll-a il
backend ogni 2s e aggiorna il grafico man mano che arrivano
nuovi sample. La pagina pausa il polling automaticamente
quando la tab del browser \`e in background, per non sprecare
banda.

\section{Lockati: lezioni che non si toccano}

\textbf{Cos'\`e.} Il flag \texttt{locked} sulla tabella
\texttt{Lesson} indica che quella lezione \emph{non deve
essere spostata} dal solver. Utile per fissare a mano una
manciata di lezioni critiche e poi lanciare l'ottimizzazione
su tutto il resto.

\textbf{Come si usa.} Dal Monitor (e dalla griglia Schedule
con click-destro), il bottone ``Blocca'' marca la lezione
come locked. Il bottone ``Sblocca'' fa l'inverso. Le lezioni
locked sopravvivono ai run successivi: vengono caricate come
\emph{vincoli aggiuntivi HARD} dal solver, che dovr\`a
incastrare il resto attorno a esse.

Il filtro ``Lockati'' nella tab del Monitor mostra solo gli
eventi che hanno almeno una lezione lockata, utile per
visualizzare velocemente cosa hai gi\`a fissato e cosa \`e
ancora libero.
